\chapter{The Deep Nature of General Relativity}
\label{chap:the-nature-of-geometry}

\section*{The Mystery}

Microscopic reality appears fundamentally wave-like. 
But as human observers, we do not inhabit an abstract, high-dimensional Hilbert space; 
we inhabit a tangible, three-dimensional geometry. 

Why do we favor the geometric framework of general relativity over the spectral domain of Hilbert space for our macroscopic reality?

As conscious observers, we are fundamentally finite informational structures. We cannot exist as diffuse, fluid, overlapping gradients.
Observers require stable subsystem boundaries that preserve their information across time. 
Based on real world experience we know that happens if two highly complex, entangled informational systems were to spatially overlap and mingle their states even
a little. Their internal data organization would instantly disrupt, causing immediate decoherence and a catastrophic loss of structural identity.
Just dip your little toe into an acid and the consequences are catastropic. 

To maintain informational boundaries within a high-dimensional, fully entangled Hilbert space is expensive, requiring an unsustainable expenditure of information.

However, boundaries can be defined very economically within a low-dimensional geometric space.
Geometry allows cheap description of locality. It establishes clear distinction between inside and outside.



\section{The Screen}

As Chapter~\nameref{chap:from-nothing-to-something} demonstrated, within the present framework, expanding spacetime and increasing entropy are two aspects of the same underlying process.

Spacetime is interpreted as the geometric representation of information compressed through a complex-valued Schrödinger codec.

Consider how conventional movies are displayed. A movie is projected onto a flat, two-dimensional screen with rectangular topology.
Its pixels have uniform sizes and distribution, regardless of how much detail a particular region contains.

The universe employs a fundamentally different arrangement.

Its geometric representation is a three-dimensional, curved manifold, with no predetermined restriction on its topology. Unlike a conventional screen, its spatial resolution need not be uniform.
Representational elements can concentrate in regions requiring greater detail while remaining sparse where less detail is sufficient.

The result is a variable-resolution geometric representation in which informational resources are distributed according to structural complexity rather than an arbitrary uniform grid.

\subsection*{Fermions: The Pixels of Reality}

A conventional display constructs images from individual pixels. In a typical RGB display, each pixel combines three elementary color components to reproduce an enormous range of colors.

What plays the role of pixels in the universe?

In the proposed interpretation, fermionic matter particles play this role. Quarks and leptons provide the localized, discrete elements through which matter and its structures become manifest in the geometric representation.

They are not necessarily tiny geometric squares. Rather, they are the localized elements of the representation that give the observer a discrete and structured physical world.

Just as a movie screen does not contain miniature versions of the objects depicted in the movie, the geometric representation need not contain miniature versions of the underlying spectral structures.
The apparent objects emerge from the organization and relationships of its representational elements.

\subsection*{Bosons: The Codec and Its Interactions}

What, then, is the role of bosons in this analogy?

Quantum field theory describes interactions between particles through the exchange of bosons. In perturbative calculations, virtual particles represent intermediate contributions to these interactions.

In the present framework, I propose a different ontological interpretation: virtual particles are not necessarily independent physical objects travelling between matter particles. They can instead be understood as elements of the underlying spectral codec, through which relationships between the localized elements are represented.

Consider two charged particles. In conventional quantum electrodynamics, their electromagnetic interaction is described through photon exchange. In the proposed interpretation, the interaction reflects the structure of the underlying spectral representation. The apparent attraction or repulsion is a manifestation of the codec's rules, rather than necessarily a literal exchange of tiny objects.

Photons provide an interesting example of the distinction.

Real photons are detectable: they can be absorbed, emitted and registered by physical detectors. However, detection always involves an interaction. A photon is registered through its effect on a detector, typically involving electronic transitions or other matter-field interactions.

We never observe a photon independently of its interactions with the rest of the physical system.

This raises a deeper question: should the exchange of virtual photons be interpreted as a literal microscopic mechanism, or as part of the mathematical encoding through which electromagnetic relationships are described?

The present framework favors the latter interpretation for virtual particles.

Real bosonic excitations, including photons, remain part of the observable physical description. Virtual particles, however, need not be assigned the same ontological status.

In this picture, fermions provide the localized elements of the geometric representation, while bosonic interactions describe how those elements relate to one another through the underlying spectral codec.


\section{Chain of Knots: the Dynamic Measuring Tape Illusion}

We are not external observers looking the pixels on the screen from outside. We are movie stars in the film
constructed from the pixels we are trying to measure. 
In such a universe, distances cannot be gauged with an absolute external metric. 
Measuring rods must be built from the emergent structures of the universe itself.

From the outside, the universe would look like a simple informational system expanding toward maximum entropy 
and flattening into a featureless state according to a parabolic saturation curve. 

But because we live inside, our rulers are made of pixels of reality. And just like a random bitstring may have sequences
of ordered information, pixels too can form ordered structures.

These structures are not added from outside; they are paid for from the fixed informational bit budget of the spacetime fabric itself.

Imagine 1D  universe (bitstring) of exactly 100 bits of white noise. Just by a change, it may contain a specific sub-string that happens to
mach the structure we label as neutron.
That neutron is not something additional information added to the universe (the $100$ bits). It is an emergent structure within the universe.
It's existence requires the informational equivalent of ten elements.

Without the neutron, the resolution of the space is $100$. With the neutron, the local spatial resolution therefore drops from $100$ to $91$. 

\begin{equation}
  R_Q(\tau)=R_{\mathrm{max}}(\tau)-\frac{1}{n}\sum\text{(bits locked in knots)}
\end{equation}

From the internal observer's perspective, whose measuring rod has itself been redefined, space appears to have contracted locally.


\subsection{Temporal resolution}

The movie of the universe consists of many frames. Just like we can ask the probability of a structure to emerge in
each spatial frame, we can ask how many frames it statistically takes the structure to emerge across frames.
Obviously, for elementary pixels (particles, bits) the probability is $100$. However, the more complex the structure, the smaller the probability.

We can model this assuming a minimal possible change between two consecutive frames; that is, one bit of information.

A structure of width \(w\) can therefore advance its own proper time only after 
a full set of \(w\) elementary flips has occurred. Its retarded clock 
is therefore the sample-and-hold quantisation

\begin{equation}
  \tau_{\mathrm{local}}(w)=w\cdot\bigl\lfloor\tau/w\bigr\rfloor,
\end{equation}

with associated lapse \(1/w\). Heavier knots age more slowly with 
respect to the universal raw coordinate \(\tau\). 

In other words, the same chain-of-knots analogy applies temporally: each knot 
introduces a coarser ticking rate, stretching the perceived duration 
of events that occur at that scale.

Emergent observers can measure ``change'' only by using the the rate at which their constituent micro-structures change.

Matter appears to curve both space and time locally. 



\section*{Geometry as Compression}

Consider an hourglass. There are an astronomical number of ways in which sand grains can arrange themselves inside its lower chamber.
Despite this enormous underlying microstate entropy, the grains tend to form a remarkably simple macroscopic structure: an approximately uniform, cone-shaped pile.

The hourglass acts as a spatial cone builder.
It maps billions of independent microscopic degrees of freedom into a macroscopic geometric structure that can be described using just a handful of variables, such as its height, base radius and angle of repose.

The microscopic arrangements differ enormously, yet their macroscopic geometric descriptions remain remarkably similar.
A small set of macroscopic constraints effectively captures the essential structure of an enormous number of microscopic configurations.

To scale this idea to its absolute limit, consider a gravitational system in a 3+1-dimensional spacetime manifold. Is there a physical upper limit to how much information can be represented within a given region?

Black holes provide a remarkable answer.

Black hole entropy is proportional to the area of the event horizon rather than the volume of the enclosed region. The Bekenstein–Hawking entropy formula establishes a fundamental relationship between geometry and information:

\[
S_{BH}=\frac{k_B A}{4l_P^2}
\]

In this sense, black holes represent an ultimate limit on the entropy that can be contained within a gravitational system of a given size.

Consider throwing arbitrary objects into a black hole: guitars, spaceships, or burning stars. Their enormously complicated external structures disappear behind the event horizon.
Once the black hole settles into a stationary state, its exterior geometry is characterized by just three macroscopic parameters: mass ($M$), electric charge ($Q$), and angular momentum ($J$).

The No-Hair Theorem captures this remarkable simplification of the observable exterior geometry. The enormous complexity of the original matter configuration is no longer reflected in the black hole's external geometric description.

This does not mean that the microscopic information has necessarily been destroyed. The relationship between black hole entropy, microscopic degrees of freedom and information preservation remains one of the deepest problems in theoretical physics.

Nevertheless, the geometric simplification is extraordinary.

So, if quantum mechanics provides a description of nature in terms of spectral structures and their associated amplitudes, what role does general relativity play?

Within the present information-theoretic framework, the proposed answer is **geometric compression**.

General relativity describes how complex physical configurations can be represented through geometric structures. From local gravitational systems to the largest cosmological scales, geometry provides a compact description of physical relationships that would otherwise require an enormously more complicated account.

The two theories can therefore be interpreted as complementary manifestations of a common principle:

Quantum mechanics provides spectral compression, while general relativity provides geometric compression.

Under this interpretation, both are optimized descriptions of different informational aspects of physical reality, unified by a single imperative:

\[
\boxed{\text{Minimization of total description length}}
\]


